\section*{2 Bounded-depth Frege systems and automatability (author source pp.22:3--22:5)}
\subsection*{2.1 Basic definitions}
We assume that formulas are built from constants 0 and 1, propositional variables and their negations, unbounded conjunctions and unbounded disjunctions. So negations can only appear next to variables. The depth of a formula is the maximum nesting of conjunctions and disjunctions in it. Formally,
\[d(0)=d(1)=d(x)=d(\neg x)=0,\qquad d(\circ\{F_1,\ldots,F_k\})=1+\max_i d(F_i),\]
where $x$ is a variable and $\circ$ is either a conjunction $\bigwedge$ or a disjunction $\bigvee$.
Depth-0 formulas that are not constants are called literals. We often write literals in the form $x^\varepsilon$, where $x^1:=x$ and $x^0:=\neg x$. Depth-1 formulas are called clauses/terms, clauses being disjunctions and terms conjunctions of literals. Depth-2 formulas that are disjunctions of terms are called DNF formulas and depth-2 formulas that are conjunctions of clauses are called CNF formulas. DNF formulas each conjunction of which consists of at most $k$ literals are called $k$-DNF formulas; $k$-CNF formulas are defined similarly. We define $\Sigma_d^{s,k}$ to be the class of all formulas $F$ for which there is a depth-$d$ formula $G$ that is semantically equivalent to $F$, the outermost connective of $G$ is $\bigvee$, and
\begin{enumerate}\item $G$ contains at most $s$ subformulas of depth at least 2;\item all depth-2 subformulas of $G$ are either $k$-DNFs or $k$-CNFs.\end{enumerate}
Similarly, $\Pi_d^{s,k}$ is defined as the class of all formulas $F$ for which there is a depth-$d$ formula $G$ semantically equivalent to $F$, the outermost connective of which is $\bigwedge$, satisfying the above two conditions.
A restriction is an assignment $\rho:V\to\bits$ of truth values to a set $V$ of variables. For a restriction $\rho$ and a formula $F$, we denote by $F|_\rho$ the formula resulting by replacing every variable $x$ of $F$ which is in the domain of $\rho$ by $\rho(x)$, and then eliminating constants from $F|_\rho$ using the identities
\[A\vee0=A,\quad A\vee1=1,\quad A\wedge0=0,\quad A\wedge1=A.\]
We call a restriction that gives a value to all variables a total assignment, or simply assignment. For a set $S$ of formulas, we write $S\models F$ if for any total assignment $\alpha$, $G|_\alpha=1$ for every $G\in S$ implies $F|_\alpha=1$. For formulas $F$ and $G$, we write $F\equiv G$ if $F$ and $G$ are semantically equivalent, i.e. it holds that $F\models G$ and $G\models F$.
\subsection*{2.2 LK proofs}
Bounded-depth Frege systems are commonly presented as subsystems of sequent calculus (LK for short) for propositional logic. We give a Tait-style formulation of LK, where we write cedents as disjunctions. The inference rules of the system are shown in Table1. There, $x$ stands for a propositional variable, $A$ and $B$ stand for arbitrary formulas whose top-most connective is $\bigvee$, $\varphi$ stands for an arbitrary propositional formula and $\Phi$ stands for a set of propositional formulas. $\bar\varphi$ is the formula that results from $\varphi$ by exchanging every occurrence of $\bigvee$ with $\bigwedge$ and vice versa, and replacing each literal $x^\varepsilon$ with $x^{1-\varepsilon}$.
\begin{center}\textbf{Table1: The rules of LK.}\par\medskip
$\begin{array}{r@{\quad}c} \text{Axioms:}&x\vee\neg x\\[8pt]\text{Weakening:}&\dfrac{A}{A\vee B}\\[12pt]\bigvee\text{-introduction:}&\dfrac{A\vee\varphi}{A\vee\bigvee\Phi},\quad\text{where }\varphi\in\Phi\\[12pt]\bigwedge\text{-introduction:}&\dfrac{A\vee\varphi_1\quad\cdots\quad A\vee\varphi_k}{A\vee\bigwedge\{\varphi_1,\ldots,\varphi_k\}}\\[12pt]\text{Cut:}&\dfrac{A\vee\varphi\quad B\vee\bar\varphi}{A\vee B}\end{array}$\end{center}
An LK proof from a set of premises $S$ is a sequence of formulas, called the lines of the proof, such that each line either belongs to $S$ or results from earlier lines by one of rules of Table1. If the last line in a proof is the empty disjunction, then the proof is called a refutation. A depth-$d$ LK proof is an LK proof each line of which is a formula of depth at most $d$. The size of a proof is the total number of symbols occurring in it.
Of particular importance among depth-$d$ LK proofs is the case of depth-1 proofs, called resolution proofs. In resolution proofs, lines are clauses, and the only applicable LK rules are the weakening and cut rule, which take the form
\[\frac C{C\vee D},\qquad\frac{C\vee x\quad D\vee\neg x}{C\vee D}\]
for clauses $C$ and $D$. In the rightmost rule, also called the resolution rule, we say that $C\vee D$ is the result of resolving $C\vee x$ on $D\vee\neg x$ on $x$.
We may view a proof as a DAG, by drawing for every line $A$, edges from the lines $A$ is derived to $A$. In case a proof DAG is a tree, we refer to the proof as being tree-like. The next propositions, due to[14], state that depth-$d$ LK proofs and tree-like depth-$(d+1)$ LK proofs can be turned into one another with only a polynomial increase in size.
\begin{originalresult}{Proposition 2.1 [14]}
A depth-$d$ LK proof of a formula $F$ from $S$ of size $s$ can be turned into a depth-$(d+1)$ tree-like LK proof of $F$ from $S$ of size polynomial in $s$.
\end{originalresult}
\begin{originalresult}{Proposition 2.2 [14,3]}
Let $S$ be a set of formulas of depth at most $d$ and $F$ a formula of depth at most $d$. A depth-$(d+1)$ tree-like LK proof of a formula $F$ from $S$ of size $s$ can be turned into a depth-$d$ LK proof of $F$ from $S$ of size $O(s^2)$.
\end{originalresult}
\subsection*{2.3 Semantic proofs, variable width and decision trees}
A semantic depth-$d$ (Frege) proof from a set of formulas $S$ is a sequence of depth-$d$ formulas $F_1,\ldots,F_t$ such that for every $i$, either $F_i\in S$ or there are $j,k<i$ such that $F_j,F_k\models F_i$. Notice that if $S$ consists of depth-$(d-1)$ formulas, then there is a trivial depth-$d$ proof of any valid consequence of $S$, as $\bigwedge S$ can be derived in $|S|-1$ steps. Thus, under this formulation, depth-$d$ proofs from $S$ are interesting only if $S$ contains depth-$d$ formulas not in $\Pi_d^{s,k}$ for any $s$ and $k$, and indeed, our results pertain to such proofs.
The definitions of lines, size of a proof, refutation, tree-like proofs, apply to semantic proofs as well. The variable width of a proof is the maximum number of variables among the lines of the proof.
Unlike size, variable width is an inherently semantic notion. In particular, it is independent of depth: any depth-$d$ proof of variable width $w$ can be transformed into a depth-1 proof of (variable) width $O(w)$. In fact, something stronger can be said. A decision tree is a binary tree the internal nodes of which are labelled by variables, and the edges by values 0 or1. Nodes query variables and the edges going from a node to its children are labelled, one by the value0 and the other by1, giving an answer to that query. No variable is repeated in a branch so that branches correspond to restrictions, and each branch has a value,0 or1, associated with it, so that the decision tree represents a Boolean function. We denote the set of branches of $T$ having the value $v$ by $\Br_v(T)$. Specifically, we say that a decision tree $T$ represents a formula $F$ if for every branch $\pi$ of $T$ with value $v$, $F|_\pi\equiv v$. The height of a decision tree is the length of its longest branch. Notice that if a formula $F$ is represented by a decision tree of height $h$, then $F\in\Sigma_2^{1,h}\cap\Pi_2^{1,h}$. We write $\height(F)$ for the minimum height of a decision tree representing $F$. The following lemma is shown in[18] for a specific type of depth-2 proofs, but holds for proofs of arbitrary depth, or for that matter, arbitrary sound proofs.
\begin{originalresult}{Lemma 2.3}
Let $S$ be a set of clauses each containing at most $h$ literals. If there is a semantic refutation of $S$ each line of which is represented by a decision tree of height at most $h$, then there is a resolution refutation of $S$ of width at most $3h$.
\end{originalresult}
\begin{proof}
Let $F_1,\ldots,F_t$ be a semantic refutation of $S$ and let $T_i$ be a decision tree of height at most $h$ representing $F_i$. We assume that $T_t$ has a single node having the value0. For a restriction $\pi$, let $C_\pi$ be the minimal clause falsified by $\pi$. We will show that for every $i$, for every branch $\pi\in\Br_0(T_i)$, we can derive $C_\pi$ via a resolution proof of width at most $3h$. Notice that $C_\pi$ for $\pi\in\Br_0(T_t)$ is the empty clause, so this construction will give a refutation.
If $F_i$ is a clause $C$ in $S$, then every $\pi\in\Br_0(T_i)$ must make every literal in $C$ false, hence $C_\pi$ is a weakening of $C$. Assume now that $F_i$ is derived from $F_j$ and $F_k$ and we have derived all clauses $C_\pi$ for $\pi\in\Br_0(T_j)\cup\Br_0(T_k)$. Let $\sigma\in\Br_0(T_i)$, and let $T$ be the tree resulting by appending a copy of $T_k$ at the end of every branch $\pi\in\Br_1(T_j)$ of $T_j$. We will use $T$ to extract a resolution proof of $C_\sigma$. More specifically, for every node $u$ of $T$ such that the path $\pi_u$ from the root of $T$ to $u$ corresponds to a restriction that is consistent with $\sigma$, we will derive $C_\sigma\vee C_{\pi_v}$. When we reach the root of $T$ we will have derived $C_\sigma$. If $u$ is a leaf of $T$, then we claim that $C_{\pi_u}$ is a weakening of some clause $C_\pi$ for $\pi\in\Br_0(T_j)\cup\Br_0(T_k)$. To see this, let $\pi_u=\pi_j\cup\pi_k$, where $\pi_j$ is the part of $\pi_u$ that belongs to $T_j$ and $\pi_k$ the part that belongs to $T_k$. Since $F_j,F_k\models F_i$ and $\pi_u$ is consistent with $\sigma$, it cannot be the case that both $\pi_j\in\Br_1(T_j)$ and $\pi_k\in\Br_1(T_k)$, otherwise a total assignment extending both $\pi_u$ and $\sigma$ would make $F_j$ and $F_k$ true, but $F_i$ false. Suppose now that $u$ is not a leaf of $T$ and suppose that $v$ and $w$ are its children. Then either $\pi_v$ and $\pi_w$ are both consistent with $\sigma$, in which case $C_\sigma\vee C_{\pi_u}$ can be derived by resolving $C_\sigma\vee C_{\pi_v}$ and $C_\sigma\vee C_{\pi_w}$ on the variable labelling $u$, or one of the children, say $v$, will be consistent with $\sigma$ and thus $C_\sigma\vee C_{\pi_u}$ will be identical to $C_\sigma\vee C_{\pi_v}$.
\end{proof}
\subsection*{2.4 Automatability and the main result}
A proof system $\sigma$ is called automatable[5] if there is an algorithm that given a set of formulas $S$ and a formula $\varphi$ provable from $S$, outputs a $\sigma$-proof of $\varphi$ from $S$ in time polynomial $r+s$, where $r$ is the total size of $S$ and $s$ the size of the shortest $\sigma$-proof of $\varphi$ from $S$.
